\documentclass[11pt]{article}
\usepackage[a4paper,margin=2.2cm]{geometry}
\usepackage{amsmath,amssymb,booktabs,microtype}
\usepackage[hidelinks]{hyperref}
\newcommand{\E}{\mathbb{E}}
\newcommand{\N}{\mathcal{N}}
\title{Evaluation of the mixture fits: recommendation, price counterfactuals and basket generation}
\author{Version-4 basket model, branch \texttt{v4-mixture}}
\date{30 October 2026}
\begin{document}
\maketitle

\noindent All numbers are on the \emph{validation weeks} (\texttt{trip\_split = 1}). These trips were
never used in training or model selection. The test weeks were not read. Tool:
\texttt{research/evaluate\_mixture.py}. Its protocols are those of the evaluators removed in
\texttt{d716f95}, so results compare with earlier ones. Only the treatment of the taste $z$ changes,
to the mixture form.

\section{What is evaluated}

\paragraph{Checkpoints.} Real data (Dunnhumby, 5,455 products):
\begin{itemize}
\item \textbf{step 360}: the old best mixture checkpoint.
\item \textbf{R2+fw}: the new recipe. It continues from step 360 to step 840 without whitening,
  then applies one final compensated whitening (\texttt{docs/training\_theory.pdf}).
\item \textbf{R0}: the old recipe continued to step 840.
\end{itemize}
Synthetic world \texttt{disc\_c48\_s3} (522 products, with causal truth for price changes):
\begin{itemize}
\item the control's best (step 420) and final (step 1,210) checkpoints;
\item S1: household and product-context blocks frozen;
\item S9: S1 with Adam's state kept across whitening, the best synthetic fit.
\end{itemize}

\paragraph{Recommendation.} One product is hidden from each basket that has at least 2 products.
Every product the store offers is then ranked by $\log p(\text{rest}+j)-\log p(\text{rest})$. Because
$Z_z$ is the same for both baskets, this equals
\[
\log \E\bigl[\exp\Delta E_z(j)\,\big|\,\text{rest}\bigr],\qquad
\Delta E_z(j)=d_j-s\cdot\Delta s_j-\tfrac12|\Delta s_j|^2+z\cdot\Delta s_j .
\]
Here $d_j$ is the exact change of the energy, and $\Delta s_j$ the change of the basket field. The
expectation is over the posterior of $z$ given the rest, with 32 importance draws. The cases are
2,000 validation trips, with a fixed seed. The table reports MRR, recall@$k$ and the median rank;
standard errors are clustered by household. ``No taste integral'' ranks by $d_j$ alone, which is
the old form's formula.

\paragraph{Generation.} One basket per validation trip: $z\sim\N(0,I)$, then $S\mid z$ exactly. A
gradient-boosted classifier tries to tell real from generated baskets, using the full basket or
only the categories present. AUC 0.5 means the two cannot be told apart. Three seeds; 6,000 trips on
real data, and all 2,755 on synthetic data.

\paragraph{Price counterfactuals.} \emph{Synthetic:} 32 scenarios, each multiplying one product's
price by 0.8 or 1.2, compared with the generator's causal truth. The quantities are the relative
change in purchases of the product, of the rest of its subcategory and of the rest of its category,
and the change in basket size. The model's incidence is the mean, over 64 common draws
$z\sim\N(0,I)$ per trip, of the exact $P_z(j\in S)$, on 400 validation trips. \emph{Real:} 166,164
validation-week discount events (a product's chain discount depth changes between weeks). The
measured within-category response comes from a conditional binomial with an offset for the
category's change, with product-clustered 95\% intervals. The model's response is the depth slope,
the display and mailer effects, and its agreement with the observed share changes. It uses 580 trips
(one per store-week) and 64 draws of $z$.

\section{Results}

\paragraph{Recommendation} (real: 1,620 cases; synthetic: 1,974).
\begin{center}\small
\begin{tabular}{llllll}
\toprule
& model & MRR & recall@10 & recall@20 & median rank\\
\midrule
real & step 360 & $0.119\pm0.007$ & 0.228 & 0.305 & 92\\
& \textbf{R2+fw} & $\mathbf{0.123\pm0.007}$ & 0.227 & 0.306 & 86\\
& R0 & $0.116\pm0.007$ & 0.227 & 0.293 & 100\\
& no taste integral (R0) & $0.076\pm0.006$ & 0.132 & 0.177 & 264\\
& popularity & $0.045\pm0.004$ & 0.075 & 0.097 & 760\\
\midrule
synthetic & control best & $0.259\pm0.008$ & 0.485 & 0.613 & 12\\
& control final & $0.277\pm0.008$ & 0.502 & 0.642 & 10\\
& S1 & $0.276\pm0.008$ & 0.509 & 0.648 & 10\\
& \textbf{S9} & $\mathbf{0.278\pm0.008}$ & 0.513 & 0.652 & 10\\
& popularity & $0.084\pm0.005$ & 0.166 & 0.280 & 50\\
\bottomrule
\end{tabular}
\end{center}
Paired differences, household-clustered:
\begin{itemize}
\item real: R2+fw $-$ step 360 $=+0.0034\pm0.0024$; R0 $-$ step 360 $=-0.0037\pm0.0025$;
\item synthetic: S9 $-$ control best $=+0.019\pm0.004$.
\end{itemize}
The model is 2.7 times better than popularity on real data and 3.3 times on synthetic data.
Integrating over the taste of the rest of the basket matters: on real data, MRR is 0.116 with it
and 0.076 with the old form's formula. On real data the differences between the three checkpoints
are small, and only the ordering (new $>$ old best $>$ old recipe continued) is consistent with
the likelihood.

\paragraph{Generation.}
\begin{center}\small
\begin{tabular}{lllllll}
\toprule
& model & full AUC & category AUC & size (real 7.50 / 8.88) & categories (5.88 / 8.11) & 20+ products\\
\midrule
real & step 360 & 0.709 & 0.661 & 8.92 & 7.42 & 0.130 (real 0.092)\\
& \textbf{R2+fw} & \textbf{0.675} & \textbf{0.633} & \textbf{7.24} & \textbf{5.96} & \textbf{0.080}\\
& R0 & 0.711 & 0.682 & 8.62 & 7.07 & 0.119\\
\midrule
synthetic & control best & 0.765 & 0.844 & 8.72 & 7.95 & \\
& control final & 0.687 & 0.788 & 7.81 & 7.19 & \\
& S1 & 0.686 & 0.786 & 8.01 & 7.37 & \\
& \textbf{S9} & \textbf{0.683} & \textbf{0.786} & 7.98 & 7.35 & \\
\bottomrule
\end{tabular}
\end{center}
On real data the new recipe gives the most realistic baskets. Its AUC falls by 0.034 (full) and
0.028 (categories) against step 360. Its size is within 3.5\% of the data, against 19\% too large
for step 360. The number of categories is within 1.4\%, against 26\% too many. On synthetic data,
the later checkpoints are much more realistic than the control's best, even though the control's
held-out likelihood fell after step 420. Their baskets are 10\% too small, the shape gap of
\texttt{training\_theory.pdf}, Section~6.

\paragraph{Price counterfactuals, synthetic} (32 scenarios; S9; the control is within 0.02 on
every number).
\begin{center}\small
\begin{tabular}{lllll}
\toprule
& correlation with truth & sign agreement & slope of model on truth & mean change at $\times0.8$ / $\times1.2$\\
\midrule
own product & 0.93 & 1.00 & 0.10 & truth $+0.47/-0.35$, model $+0.05/-0.04$\\
rest of subcategory & 0.83 & 0.97 & 0.06 & \\
rest of category & 0.92 & 0.95 & 0.13 & \\
basket size & $-0.19$ & 0.56 & & truth and model both $\approx0$\\
\bottomrule
\end{tabular}
\end{center}
\textbf{The directions and the ranking of the scenarios are right; the sizes are about ten times
too small.} A 20\% price cut raises the product's purchases by 47\% in the truth and by 5\% in the
model, so the own-price elasticity is about $-2$ in the truth and $-0.25$ in the model. Substitution
to the rest of the subcategory and the category is also 6--13\% of the truth. This holds for every
synthetic checkpoint. In S1 and S9 the price parameters ($\gamma_h$, $\beta$) were not trained by
the mixture fit, and the control is no better, so the gap comes from the base model's price terms,
not from the mixture training.

\paragraph{Price and promotion response, real} (validation weeks; all three checkpoints within
0.005 of each other, because the real fits did not train the price and promotion terms).
\begin{center}\small
\begin{tabular}{lll}
\toprule
& measured (product-clustered 95\%) & model\\
\midrule
discount depth, within-category share slope & $0.64\ [0.35,\ 0.93]$ & 1.38 (whole week) -- 1.50 (own price only)\\
display on vs off, share log ratio & $0.375\ [0.28,\ 0.47]$ & 0.37 (purchase-weighted), 0.52 (mean)\\
mailer on vs off & $0.347\ [0.26,\ 0.43]$ & 0.39 (purchase-weighted), 0.51 (mean)\\
\midrule
agreement with observed share changes & & correlation 0.16, sign agreement 0.61\\
mean share change when the discount deepens & $+0.108$ & $+0.015$\\
\bottomrule
\end{tabular}
\end{center}
Display and mailer effects match the measured ones, weighted by purchases. The response to the
discount depth is about twice the measured one, and outside its interval. Agreement event by event
is weak, as before (\texttt{v4\_comp\_promo}: 1.38--1.51, correlation 0.15--0.16, the same).

\section{Verdict}

\begin{enumerate}
\item \textbf{Recommendation:} good, and the new recipe is the best on real data. The model is
  2.7--3.3 times better than popularity, and the taste integral is worth 50\% of the MRR on real
  data.
\item \textbf{Basket generation:} the new recipe is clearly the most realistic on real data:
  lowest classifier AUC, size within 3.5\%, categories within 1.4\%. A classifier can still tell
  real from generated baskets (AUC 0.68 full, 0.63 categories).
\item \textbf{Price counterfactuals: not good enough yet.} On synthetic data the signs and the
  ranking are right, but the sizes are about 10 times too small. On real data the discount
  response is twice the measured one, while display and mailer match. In neither case does the
  mixture training cause this. It sits in the price terms of the base model, which the
  mixture fits keep fixed (real) or near their start (synthetic). This is the next problem to
  work on. The goal requires correct price counterfactuals, and none of today's fixes touched
  them.
\end{enumerate}

\end{document}
